% 2009 年湖北卷文科第 10 题：三角形数与正方形数的点阵图（图 1、图 2）。
% 来源 PDF：普通高考/2009/2009湖北文.pdf，SHA-256 f60ba7a5a2fef3449fba9474fa973c616aa52a15926b3a97f0dda6b909f4c861。
% 原页：第 1 页第 2 栏中部，书页码 525。
% 重绘依据：按原页逐点重绘，点数与排列照原页。
%   图 1：三角形数 1、3、6、10，自下而上每行比下一行少一点且逐行居中（1+2+3+4 点）。
%   图 2：正方形数 1、4、9、16，为 1×1、2×2、3×3、4×4 的方阵，各组底边对齐。
%   两组末均以省略号表示后续。
\begin{tikzpicture}[x=0.46cm, y=0.46cm, semithick,
  every node/.style={inner sep=1pt, font=\small}]

  % ---------- 图 1：三角形数 ----------
  % 1
  \fill (0,0) circle (2.1pt);
  % 3
  \fill (2.4,0) circle (2.1pt);
  \fill (3.4,0) circle (2.1pt);
  \fill (2.9,1) circle (2.1pt);
  % 6
  \fill (5.8,0) circle (2.1pt);
  \fill (6.8,0) circle (2.1pt);
  \fill (7.8,0) circle (2.1pt);
  \fill (6.3,1) circle (2.1pt);
  \fill (7.3,1) circle (2.1pt);
  \fill (6.8,2) circle (2.1pt);
  % 10
  \fill (10.2,0) circle (2.1pt);
  \fill (11.2,0) circle (2.1pt);
  \fill (12.2,0) circle (2.1pt);
  \fill (13.2,0) circle (2.1pt);
  \fill (10.7,1) circle (2.1pt);
  \fill (11.7,1) circle (2.1pt);
  \fill (12.7,1) circle (2.1pt);
  \fill (11.2,2) circle (2.1pt);
  \fill (12.2,2) circle (2.1pt);
  \fill (11.7,3) circle (2.1pt);

  \node at (14.7,1.4) {$\cdots$};
  \node[below] at (0,-0.6) {$1$};
  \node[below] at (2.9,-0.6) {$3$};
  \node[below] at (6.8,-0.6) {$6$};
  \node[below] at (11.7,-0.6) {$10$};
  \node at (7.3,-2.55) {图 1};

  % ---------- 图 2：正方形数 ----------
  \begin{scope}[yshift=-3.5cm]
    % 1
    \fill (0,0) circle (2.1pt);
    % 4
    \fill (2.4,0) circle (2.1pt);
    \fill (3.4,0) circle (2.1pt);
    \fill (2.4,1) circle (2.1pt);
    \fill (3.4,1) circle (2.1pt);
    % 9
    \foreach \i in {5.8,6.8,7.8} {
      \foreach \j in {0,1,2} {
        \fill (\i,\j) circle (2.1pt);
      }
    }
    % 16
    \foreach \i in {10.2,11.2,12.2,13.2} {
      \foreach \j in {0,1,2,3} {
        \fill (\i,\j) circle (2.1pt);
      }
    }

    \node at (14.7,1.4) {$\cdots$};
    \node[below] at (0,-0.6) {$1$};
    \node[below] at (2.9,-0.6) {$4$};
    \node[below] at (6.8,-0.6) {$9$};
    \node[below] at (11.7,-0.6) {$16$};
    \node at (7.3,-2.55) {图 2};
  \end{scope}
\end{tikzpicture}
